% Procedencia: COMPATIBILIDAD_HOLONOMIA_E_INFORMACION.md, secciones 2--9.
% Desarrollo científico conservado; gestión y correspondencia en archivo separado.
\section{Compatibilidad constitutiva, holonomía relativa e información operacional}
\label{md:comp:capitulo}

\subsection{Reciprocidad diferencial de velocidad e impedancia}
\label{md:comp:reciprocidad}

El antecedente constitutivo de la sección~\ref{iii:sec:constitutiva} construye, sobre las hojas
conservadas de su representación,
\begin{equation}
 \begin{gathered}
 T=e^{-xI-y\mathcal R},\qquad
 r_\pm=f\bigl(e^{-30(x\pm y)}\bigr),\\
 V=(I-T^{90})(I-T^{120})^{-1}=r_+P_++r_-P_-.
 \end{gathered}
 \label{md:comp:constitutiva}
\end{equation}
donde
\[
 f(s)=\frac{1-s^3}{1-s^4},\qquad x>|y|.
\]
Los proyectores de hoja permanecen fijados. La cámara física marcada
\(x>y>0\) y su conjugada están contenidas en esta colección analítica.
Definimos las coordenadas logarítmicas normalizadas
\begin{equation}
 c_\ell=\log\widehat c=-\log r_+-\log r_-,\qquad
 z_\ell=\log\widehat Z=\log r_--\log r_+.
 \label{md:comp:logaritmos}
\end{equation}
Estas letras no designan nuevas constantes ni el contraángulo \(C^*\).
Sean
\[
 u=30(x+y),\qquad v=30(x-y),\qquad
 g(w)=\log f(e^{-w}),\qquad a=g'(u),\quad b=g'(v).
\]

\begin{proposition}[Reciprocidad constitutiva diferencial]
\label{md:comp:prop-reciprocidad}
En el dominio \(x>|y|\), se tiene
\begin{equation}
 g'(w)=\frac{3}{e^{3w}-1}-\frac{4}{e^{4w}-1}>0,
 \qquad
 D(c_\ell,z_\ell)=-30
 \begin{pmatrix}a+b&a-b\\a-b&a+b\end{pmatrix}.
 \label{md:comp:jacobiano}
\end{equation}
Los autovalores del jacobiano son \(-60a\) y \(-60b\). Por tanto, es
simétrico y definido negativo, su determinante es \(3600ab>0\), y
\begin{equation}
 \boxed{\partial_y\log\widehat c=\partial_x\log\widehat Z.}
 \label{md:comp:reciprocidad-cruzada}
\end{equation}
\end{proposition}
\begin{proof}
Se derivan las identidades
\(c_\ell=-g(u)-g(v)\) y \(z_\ell=g(v)-g(u)\).
La positividad de \(g'\) procede del decrecimiento estricto, en \(k>0\),
de \(k/(e^{kw}-1)\). También resulta de la expresión racional positiva
de la sección~\ref{md:comp:restitucion}. La matriz obtenida tiene los
autovalores y el determinante indicados; la igualdad de sus entradas
no diagonales da \eqref{md:comp:reciprocidad-cruzada}.
\end{proof}

La igualdad relaciona la respuesta de la velocidad a la separación
orientada con la respuesta de la impedancia a la coordenada común.
Ambas sensibilidades proceden de una misma función espectral. Es una
reciprocidad de esta colección constitutiva; no se identifica sin prueba
con una ley termodinámica o de electromagnetismo macroscópico.

\subsection{Potencial espectral y conexión de profundidad dos con Catalán}
\label{md:comp:potencial}

\subsubsection{Construcción del potencial}
\label{md:comp:construccion-potencial}

Desde los canales ya construidos, definimos
\begin{equation}
 \mathscr F(x,y)=\sum_{\sigma=\pm1}\sum_{n\geq1}
 \left[
 \frac{e^{-120n(x+\sigma y)}}{120n^2}
 -\frac{e^{-90n(x+\sigma y)}}{90n^2}
 \right].
 \label{md:comp:serie-potencial}
\end{equation}

\begin{proposition}[Potencial de las dos respuestas logarítmicas]
\label{md:comp:prop-potencial}
La serie \eqref{md:comp:serie-potencial} define en \(x>|y|\) un potencial
que satisface
\begin{equation}
 \boxed{\partial_x\mathscr F=c_\ell,\qquad
        \partial_y\mathscr F=z_\ell.}
 \label{md:comp:gradiente-potencial}
\end{equation}
Su hessiano es el jacobiano de
\eqref{md:comp:jacobiano}; en consecuencia, \(\mathscr F\) es
estrictamente cóncava. La constante aditiva queda fijada por
\(\mathscr F\longrightarrow0\) cuando \(x-|y|\longrightarrow\infty\).
\end{proposition}
\begin{proof}
En todo compacto del dominio \(x>|y|\), las cantidades \(x\pm y\)
tienen un mínimo positivo. Una potencia de \(n\) multiplicada por una
exponencial decreciente es sumable; por ello, la serie y sus derivadas
de cualquier orden fijo convergen uniformemente en ese compacto.
Derivando término a término y utilizando
\(-\log(1-z)=\sum_{n\geq1}z^n/n\), se obtienen las dos identidades
de \eqref{md:comp:gradiente-potencial}. El hessiano coincide entonces
con \eqref{md:comp:jacobiano}, que es definido negativo.
La normalización en el infinito es la de la serie exhibida.
\end{proof}

\(\mathscr F\) es un potencial matemático de la respuesta. No se le
asignan unidades de energía ni se declara un potencial termodinámico.

\subsubsection{Realización por un mismo momento espectral}
\label{md:comp:momento-espectral}

Sea
\[
 D_2(z)=\sum_{n\geq1}\frac{z^n}{n^2}.
\]
Su reconocimiento como dilogaritmo es posterior a esta definición por
serie. Para \(|z|<1\), satisface
\((z\partial_z)^2D_2(z)=z/(1-z)\). En esta notación,
\begin{equation}
 \mathscr F=
 \sum_{\sigma\in\{+,-\}}
 \left[\frac{D_2(q_\sigma^{120})}{120}
       -\frac{D_2(q_\sigma^{90})}{90}\right].
 \label{md:comp:potencial-dos-momentos}
\end{equation}
El momento orientado de Catalán construido en la sección~\ref{iii:sec:catalan} satisface
\begin{equation}
 \mathcal C_2(s)=\operatorname{Im}D_2(is)
 =\sum_{k\geq0}\frac{(-1)^ks^{2k+1}}{(2k+1)^2},\qquad
 \mathcal C_2(1)=G_{\rm Cat}.
 \label{md:comp:catalan-momento}
\end{equation}

\begin{proof}
Las potencias pares de \(i\) son reales y las impares tienen parte
imaginaria \((-1)^k\). Esto da la serie de
\eqref{md:comp:catalan-momento}, que converge absolutamente incluso
en \(s=1\). Para \(0\leq s<1\), aplicar término a término dos veces
\(s\partial_s\) devuelve \(s/(1+s^2)\), el coeficiente orientado del
resolvente del cuarto de giro ya presente en el corpus.
La identidad diferenciada en \(s=1\), si se utiliza, se interpreta
por continuación analítica o límite de Abel: la serie dos veces
diferenciada no converge ordinariamente en ese extremo.
La identificación \(\mathcal C_2(1)=G_{\rm Cat}\) procede directamente
de la serie original absolutamente convergente.
\end{proof}

Por tanto, el momento de Catalán y el potencial constitutivo se reúnen
mediante un mismo momento espectral de profundidad dos, evaluado en
soportes diferentes: el cuarto de giro orientado y las contracciones
de los sectores \(90/120\). Esta composición no identifica Catalán
con \(\mathscr F\), ni afirma que su único valor extremo determine
todos los canales. La colección analítica y sus argumentos conservan
información que un valor aislado no contiene.

\subsubsection{Compatibilidad con amplificación}
\label{md:comp:amplificacion-potencial}

Para \(T_m=T\otimes I_{9^m}\) y
\(\tau_m=\operatorname{Tr}/(2\cdot9^m)\), se tiene exactamente
\begin{equation}
 \mathscr F_m=2\tau_m\left[
 \frac{D_2(T_m^{120})}{120}
 -\frac{D_2(T_m^{90})}{90}\right]=\mathscr F.
 \label{md:comp:potencial-amplificado}
\end{equation}
La multiplicidad \(9^m\) cancela la normalización. El mismo argumento
conserva las derivadas. Esto acredita la amplificación de la
representación, no cualquier refinamiento imaginable del TPK.
No se sustituye el operador temporal por una lista de autovalores
sin marcas de hoja.

Para cualquier lazo suave cerrado en esta colección analítica,
\begin{equation}
 \oint(c_\ell\,dx+z_\ell\,dy)=0.
 \label{md:comp:integral-cerrada}
\end{equation}
Esta igualdad se refiere a la proyección constitutiva de dos
coordenadas. No implica que la historia enriquecida o su holonomía
operatoria sean triviales.

\subsection{Restitución exacta y sensibilidad finita}
\label{md:comp:restitucion}

\begin{proposition}[Restitución analítica de las coordenadas angulares]
\label{md:comp:prop-restitucion}
Con \(L=\log(16/9)\), el mapa
\((x,y)\longmapsto(c_\ell,z_\ell)\) es un difeomorfismo real analítico
del dominio \(x>|y|\) sobre
\begin{equation}
 0<c_\ell+z_\ell<L,\qquad 0<c_\ell-z_\ell<L.
 \label{md:comp:imagen-logaritmica}
\end{equation}
\end{proposition}
\begin{proof}
La función \(g\) crece estrictamente de \(\log(3/4)\) a \(0\).
Las combinaciones
\[
 c_\ell+z_\ell=-2g(u),\qquad c_\ell-z_\ell=-2g(v)
\]
recuperan \(u,v\) mediante \(g^{-1}\). Entonces
\(x=(u+v)/60\) y \(y=(u-v)/60\).
La inyectividad se mantiene al restringir a la imagen efectivamente
producida por HMT; no se declara que HMT produzca todo el cono.
\end{proof}

Con \(s=e^{-w}\), la derivada tiene la forma racional
\begin{equation}
 a(w)=g'(w)=
 \frac{s^3(s^2+2s+3)}{(1+s+s^2)(1+s+s^2+s^3)},
 \qquad 0<a(w)<\frac12.
 \label{md:comp:derivada-racional}
\end{equation}
La función es estrictamente decreciente, con \(a(0+)=1/2\), y
\(a(w)\sim3e^{-3w}\) en el infinito. Para probar el decrecimiento,
se utiliza
\begin{equation}
 a'(w)=-\frac{9}{4\sinh^2(3w/2)}
        +\frac{16}{4\sinh^2(2w)}<0,
 \label{md:comp:decrecimiento}
\end{equation}
pues \(k/\sinh(kw/2)\) decrece estrictamente en \(k>0\).
Además,
\begin{equation}
 \frac12e^{-3w}<a(w)<3e^{-3w}.
 \label{md:comp:cotas-exponenciales}
\end{equation}
La cota inferior resulta al multiplicar los denominadores y utilizar
\begin{equation}
 \begin{split}
 &2(s^2+2s+3)-(1+s+s^2)(1+s+s^2+s^3)\\
 &\hspace{1cm}=(1-s)(5+7s+6s^2+3s^3+s^4)>0.
 \end{split}
 \label{md:comp:polinomio-cota}
\end{equation}
La superior resulta de
\[
 3(1+s+s^2)(1+s+s^2+s^3)-(s^2+2s+3)>0.
\]

Si \(m=30(x-|y|)\) y \(M=30(x+|y|)\), la norma del inverso del
jacobiano es \(1/[60a(M)]\). Para no confundir el mapa vectorial con
el potencial escalar, escribimos
\(\mathbf F=(c_\ell,z_\ell)\).

\begin{proposition}[Cotas finitas de restitución]
\label{md:comp:prop-sensibilidad}
Sobre una región convexa en la que ambos argumentos \(u,v\) pertenecen
a \([m_0,M_0]\subset(0,\infty)\), se verifican
\begin{equation}
 60a(M_0)\|p-q\|
 \leq\|\mathbf F(p)-\mathbf F(q)\|
 \leq60a(m_0)\|p-q\|.
 \label{md:comp:cotas-finitas}
\end{equation}
\end{proposition}
\begin{proof}
Se integra el hessiano de \(\mathscr F\), es decir, el jacobiano
de \(\mathbf F\), sobre el segmento. Para la cota inferior se utiliza
\[
 -\langle\mathbf F(p)-\mathbf F(q),p-q\rangle
 \geq60a(M_0)\|p-q\|^2
\]
y la desigualdad de Cauchy--Schwarz. La cota superior utiliza la
norma máxima del jacobiano. Son cotas de errores finitos en el
dominio especificado.
\end{proof}

En \(y=0\) y cuando \(x\longrightarrow\infty\), ambas direcciones
se contraen por igual: el número de condición matricial es uno,
pero la norma del inverso crece como \(e^{90x}/180\).
Invertibilidad exacta y recuperación robusta frente a error absoluto
son, por tanto, propiedades diferentes. Cuando \(|y|\longrightarrow x\)
con \(x>0\) fijo, no hay singularidad diferencial: una derivada tiende
a \(1/2\) y la otra a \(a(60x)>0\).

\subsection{El cuarto de giro procede del ciclo dodecafásico}
\label{md:comp:cuarto-giro}

Se recupera la construcción de \eqref{iii:eq:quarter-intertwiner}. En \(\mathbb R^{12}\),
sean
\[
 C_{12}e_j=e_{j+1\bmod12},\qquad
 \Pi_4=P_4^{\rm cyc}\frac{I-C_{12}^6}{2},
\]
y
\begin{equation}
 \begin{gathered}
 u_0=(1,0,-1,0,1,0,-1,0,1,0,-1,0)^{\mathsf T},\qquad
 v_0=C_{12}u_0,\\
 W=(u_0,v_0)/\sqrt6.
 \end{gathered}
 \label{md:comp:isometria-ciclica}
\end{equation}
El símbolo \(\Pi_4\) distingue el proyector de las coordenadas de
acción \(Q,P\). Se verifican
\begin{equation}
 \begin{gathered}
 W^{\mathsf T}W=I,\qquad WW^{\mathsf T}=\Pi_4,\qquad
 C_{12}W=WJ_0,\\
 J_0=\begin{pmatrix}0&-1\\1&0\end{pmatrix},\qquad J_0^2=-I.
 \end{gathered}
 \label{md:comp:entrelazador-ciclico}
\end{equation}
En particular, \(C_{12}^3W=-WJ_0\). El plano tiene así un entrelazador
explícito con el ciclo; no se ha introducido un giro externo para
ajustar una respuesta.

El antecedente radial de \eqref{rec:eq:entrelazamiento} realiza
\[
 S_\zeta=\operatorname{diag}(e^{\zeta/2},e^{-\zeta/2}),
 \qquad J_\zeta=S_\zeta J_0S_\zeta^{-1}.
\]
Sobre ese plano definimos
\begin{equation}
 U_\zeta(\theta)=S_\zeta e^{-\theta J_0}S_\zeta^{-1}
 =\begin{pmatrix}
 \cos\theta&e^\zeta\sin\theta\\
 -e^{-\zeta}\sin\theta&\cos\theta
 \end{pmatrix}.
 \label{md:comp:transporte-plano}
\end{equation}
El signo de \(\theta\) fija la orientación elegida. La razón generada
\(\xi=C^*/A\), con \(|\xi|<1\) y ambos ángulos expresados en la misma
unidad, satisface \(\zeta=\operatorname{artanh}\xi\).
En la orientación conjugada,
\((\varepsilon,\xi)\mapsto(-\varepsilon,-\xi)\) conserva la sección
positiva de acción mediante el recuperador orientado del corpus;
\(\zeta\) cambia de signo.

\subsection{Lazo relativo de dos transportes}
\label{md:comp:lazo}

Sean
\begin{equation}
 U=U_{\zeta_s}(\theta),\qquad
 V=U_{\zeta_r}(\vartheta),\qquad
 H_{\rm lazo}=UVU^{-1}V^{-1}.
 \label{md:comp:def-lazo}
\end{equation}
Los productos se leen con acción de derecha a izquierda.
La letra \(\vartheta\) designa aquí la segunda fase, para distinguirla
de la coordenada regional \(\varphi\). Pongamos
\[
 \delta=\zeta_s-\zeta_r,\qquad
 \chi=\sinh\delta\,\sin\theta\,\sin\vartheta.
\]

\begin{proposition}[Traza y espectro del lazo relativo]
\label{md:comp:prop-lazo}
La composición \eqref{md:comp:def-lazo} satisface
\begin{equation}
 \boxed{\begin{gathered}
 \det H_{\rm lazo}=1,\qquad
 \operatorname{tr}H_{\rm lazo}=2+4\chi^2,\\
 H_{\rm lazo}=I\ \Longleftrightarrow\ \chi=0.
 \end{gathered}}
 \label{md:comp:traza-lazo}
\end{equation}
Sus autovalores son
\begin{equation}
 \lambda_\pm=(\sqrt{1+\chi^2}\pm|\chi|)^2
 =\exp\bigl(\pm2\operatorname{arsinh}|\chi|\bigr).
 \label{md:comp:espectro-lazo}
\end{equation}
\end{proposition}
\begin{proof}
Conjugando simultáneamente por \(S_{\zeta_r}^{-1}\), el operador \(V\)
es \(R(-\vartheta)\), y \(U\) tiene términos
\(e^{\pm\delta}\sin\theta\). La multiplicación directa da, en esa
referencia,
\[
 UV-VU=-2\chi\operatorname{diag}(1,-1).
\]
Como
\[
 H_{\rm lazo}-I=(UV-VU)U^{-1}V^{-1},
\]
se obtiene la equivalencia con \(\chi=0\). La multiplicación, o la
identidad de trazas para matrices unimodulares, da
\(\operatorname{tr}H_{\rm lazo}=2+4\chi^2\).
El polinomio característico es
\(\lambda^2-(2+4\chi^2)\lambda+1\), y sus raíces son las de
\eqref{md:comp:espectro-lazo}.
\end{proof}

Esto es una holonomía relativa de la composición declarada.
No se identifica con \(\Gamma_9\) ni con una actualización completa
del TPK por semejanza de nombres.

\subsubsection{Especialización enteramente determinada por el par conjugado}
\label{md:comp:lazo-conjugado}

Sean \(U_+=U_\zeta(\pi/2)\) y
\(U_-=U_{-\zeta}(\pi/2)\). Ambos satisfacen
\(U_+^2=U_-^2=-I\), pero
\begin{equation}
 \begin{aligned}
 U_+U_-&=-\operatorname{diag}(e^{2\zeta},e^{-2\zeta}),\\
 H_{\rm lazo}=(U_+U_-)^2
 &=\operatorname{diag}(e^{4\zeta},e^{-4\zeta}).
 \end{aligned}
 \label{md:comp:lazo-diagonal}
\end{equation}
Al escribir \(\xi=\tanh\zeta=C^*/A\), con ambos ángulos expresados
en la misma unidad, se obtiene
\begin{equation}
 \begin{aligned}
 H_{\rm lazo}
 &=\operatorname{diag}\left(
 \left[\frac{1+\xi}{1-\xi}\right]^2,
 \left[\frac{1-\xi}{1+\xi}\right]^2\right),\\
 \operatorname{tr}H_{\rm lazo}-2
 &=\frac{16\xi^2}{(1-\xi^2)^2}.
 \end{aligned}
 \label{md:comp:lazo-razon-angular}
\end{equation}
La fórmula de \(H_{\rm lazo}\) compone el cuarto de giro y las formas
conjugadas del mismo par. Su expresión es exacta en la razón angular
generada \(\xi\).

En la rama positiva de acción del corpus,
\(A_{\deg}=1000\alpha\) y
\(C^*_{\deg}=2(\eta_{\rm ret}+\alpha)\), con
\(\hbar_{\rm ret}=s_0\eta_{\rm ret}\). Por tanto,
\begin{equation}
 \begin{aligned}
 \xi&=\frac{\eta_{\rm ret}+\alpha}{500\alpha},\\
 \rho_{\rm lazo}=e^{4\zeta}
 &=\left(\frac{501\alpha+\eta_{\rm ret}}
                 {499\alpha-\eta_{\rm ret}}\right)^2,
 \qquad 0<\eta_{\rm ret}<499\alpha.
 \end{aligned}
 \label{md:comp:multiplicador-lazo}
\end{equation}
Los coeficientes \(501\) y \(499\) son \(500\pm1\);
\(500\) procede del factor \(1000/2\) del par angular y no es un
nuevo parámetro. Con los ejes orientados conservados, la restitución es
\begin{equation}
 \boxed{
 \eta_{\rm ret}=\alpha\left[
 500\frac{\sqrt{\rho_{\rm lazo}}-1}
          {\sqrt{\rho_{\rm lazo}}+1}-1\right].}
 \label{md:comp:restitucion-accion}
\end{equation}
La derivada
\[
 \frac{d\eta_{\rm ret}}{d\rho_{\rm lazo}}
 =\frac{500\alpha}
 {\sqrt{\rho_{\rm lazo}}(\sqrt{\rho_{\rm lazo}}+1)^2}
\]
es positiva. La cámara \(\eta_{\rm ret}>0\) corresponde a
\(\rho_{\rm lazo}>(501/499)^2\); el borde
\(\eta_{\rm ret}\longrightarrow499\alpha\) corresponde a
\(\rho_{\rm lazo}\longrightarrow\infty\).
Es una lectura inversa del coeficiente adimensional de acción desde
el transporte relativo y la coordenada \(\alpha\) ya construida.
Conserva la escala \(s_0\) como antecedente: una transformación
unimodular de forma no determina por sí sola esa escala dimensional.
La orientación conjugada invierte \(\rho_{\rm lazo}\) y transporta
la etiqueta de acción; no convierte \(\eta_{\rm ret}\) en acción negativa.

\subsubsection{Iteración del lazo y acumulación orientada}
\label{md:comp:iteracion}

Sea
\[
 \rho_{\rm lazo}=e^{4\zeta}
 =\left(\frac{1+\xi}{1-\xi}\right)^2.
\]
Para todo entero \(n\),
\begin{equation}
 \begin{gathered}
 H_{\rm lazo}^n
 =\operatorname{diag}(\rho_{\rm lazo}^{n},\rho_{\rm lazo}^{-n}),
 \\
 Q_n=\rho_{\rm lazo}^{n}Q_0,\quad
 P_n=\rho_{\rm lazo}^{-n}P_0.
 \end{gathered}
 \label{md:comp:iteracion-diagonal}
\end{equation}
La prueba es la multiplicación de matrices diagonales, extendida a
\(n<0\) mediante el inverso. Se conservan el producto
\(Q_nP_n=Q_0P_0\) y el área simpléctica, mientras el estado cambia
para \(\zeta\ne0\) y \(X_0\ne0\). En el programa de cuatro
operaciones, los incrementos de fase declarados suman cero en cada
repetición; esa contabilidad de fases no determina el transporte
total del estado.

Para \(\zeta\ne0\) y \(Q_0\ne0\), se recupera exactamente
\[
 n=\frac{\log|Q_n/Q_0|}{4\zeta}.
\]
Si \(Q_0=0\), la condición \(P_0\ne0\) permite usar
\[
 n=-\frac{\log|P_n/P_0|}{4\zeta}.
\]
La lectura requiere conocer la preparación y conservar una referencia.
La extensión multiplicativa
\(H_{\rm lazo}^{n+m}=H_{\rm lazo}^{n}H_{\rm lazo}^{m}\)
corresponde a un incremento logarítmico aditivo \(4(n+m)\zeta\).
Invertir el lazo cambia el signo de ese incremento.

Así, un programa fijo y repetido admite una respuesta que registra
su iteración. Esto es una realización matemática de acumulación
operacional en la composición estudiada. No se identifica \(n\)
con todo el acarreo o la memoria de \(\Gamma_9\), ni se deduce una
entropía nula para la dinámica completa a partir de la periodicidad
de su programa de fases. Si la amplitud inicial es desconocida y
sólo se dispone del estado final, \(n\) no queda determinado por
esa lectura aislada.

\subsubsection{Levantamiento explícito a las doce fases}
\label{md:comp:levantamiento}

Sean
\[
 \widetilde S_\zeta=I-\Pi_4+WS_\zeta W^{\mathsf T},\qquad
 L_\pm=\widetilde S_{\pm\zeta}C_{12}^3
       \widetilde S_{\pm\zeta}^{-1}.
\]
En \(\operatorname{Ran}\Pi_4\) actúan \(U_+,U_-\), respectivamente;
en el complemento, ambos actúan como \(C_{12}^3\). Por tanto,
\begin{equation}
 L_+L_-L_+^{-1}L_-^{-1}
 =I-\Pi_4+WH_{\rm lazo}W^{\mathsf T}.
 \label{md:comp:levantamiento-lazo}
\end{equation}
La acción complementaria se cancela. La traza del operador de doce
fases es \(10+\operatorname{tr}H_{\rm lazo}\).
El entrelazamiento y la composición son exactos y están comprobados
racionalmente usando los vectores enteros \(u_0,v_0\).
Este levantamiento reúne el antecedente dodecafásico y el transporte
elíptico; no afirma que su secuencia sea por sí misma la actualización
autónoma original del TPK.

\subsubsection{Refinamiento y límite}
\label{md:comp:limite}

La amplificación
\(H_{{\rm lazo},m}=H_{\rm lazo}\otimes I_{9^m}\) conserva
\[
 \frac{1}{9^m}\operatorname{Tr}H_{{\rm lazo},m}
 =\operatorname{Tr}H_{\rm lazo},
\]
así como la norma y el espectro, salvo las multiplicidades.
Con las inclusiones \(v\mapsto v\otimes e_0\), se cumple
\[
 H_{{\rm lazo},m+1}\iota_m=\iota_m H_{{\rm lazo},m}.
\]
El operador uniforme define en el límite inductivo hilbertiano una
extensión acotada, con la misma norma. Es el límite de esta
amplificación; no se sustituye por él otro sistema de refinamientos
del corpus.

\subsection{Qué distingue un cambio de coordenatización de un efecto relativo}
\label{md:comp:carta-efecto}

El transporte de coordenadas
\[
 T_n=S_{\zeta_{n+1}}R(-\theta_n)S_{\zeta_n}^{-1}
\]
telescopa:
\begin{equation}
 T_{N-1}\cdots T_0
 =S_{\zeta_N}R\left(-\sum_{n=0}^{N-1}\theta_n\right)
 S_{\zeta_0}^{-1}.
 \label{md:comp:telescopaje}
\end{equation}
Si retornan coordenatización y fase, resulta \(I\). El conmutador de la
sección~\ref{md:comp:lazo} compone evoluciones de formas distintas
en una referencia común; no es ese cambio pasivo.

Bajo un cambio común de representación \(P\), los operadores
\(U,V,H_{\rm lazo}\) se conjugan por \(P\).
La traza y el espectro de \(H_{\rm lazo}\) permanecen invariantes.
Si además \(X\mapsto PX\) y
\(g\mapsto P^{-\mathsf T}gP^{-1}\), se conserva
\(X^{\mathsf T}gX\) y toda lectura energética correctamente
transportada. Así, un \(H_{\rm lazo}\ne I\) no desaparece por
renombrar coordenadas.

Esto constituye un protocolo matemático relacional cerrado.
Su implementación experimental requiere que las evoluciones e
inversas sean operaciones físicas disponibles y que sistema,
referencia y controlador tengan una realización común.
No se afirma aquí una fuerza nueva ni una fuente de energía.

\subsubsection{Conservación de área y conservación de una misma energía}
\label{md:comp:area-energia}

\begin{proposition}[Ausencia de una métrica positiva común]
\label{md:comp:prop-metrica}
Cada cuarto de giro conserva su propia métrica:
\(g_\zeta=\operatorname{diag}(e^{-\zeta},e^\zeta)\) para \(U_+\),
y \(g_{-\zeta}\) para \(U_-\). Cuando \(\zeta\ne0\), no existe
una métrica positiva común.
\end{proposition}
\begin{proof}
Las condiciones
\[
 G=\begin{pmatrix}a&b\\b&d\end{pmatrix}>0,
 \qquad U_+^{\mathsf T}GU_+=G
\]
obligan a \(b=0\) y \(d=ae^{2\zeta}\).
La condición \(U_-^{\mathsf T}GU_-=G\) exige
\(d=ae^{-2\zeta}\). Ambas son compatibles sólo con \(\zeta=0\).
\end{proof}

Por tanto, conservar el área simpléctica no equivale a conservar
una misma forma energética durante toda composición de estos
transportes. En la referencia
\(g_r=S_r^{-\mathsf T}S_r^{-1}\), con \(\omega_r>0\),
\(X\ne0\) y \(E_r(X)=\omega_r X^{\mathsf T}g_rX/2\), se tiene
\begin{equation}
 \begin{gathered}
 \frac{E_r(H_{\rm lazo}X)}{E_r(X)}
 =\frac{x^{\mathsf T}\overline{H}_{\rm lazo}^{\mathsf T}
                   \overline{H}_{\rm lazo}x}{x^{\mathsf T}x},\\
 \overline{H}_{\rm lazo}=S_r^{-1}H_{\rm lazo}S_r,\qquad
 x=S_r^{-1}X.
 \end{gathered}
 \label{md:comp:cociente-energia}
\end{equation}
Los extremos son los cuadrados de los valores singulares.
En el caso conjugado diagonal son \(e^{\pm8|\zeta|}\), mientras
las dilataciones del estado son \(e^{\pm4\zeta}\).
Esta diferencia de exponentes es necesaria. Cambiar físicamente
de operación exige contabilizar el controlador; la ganancia de
una lectura no acredita creación de energía.

\subsubsection{El orden requiere una lectura orientada}
\label{md:comp:orden-orientado}

Intercambiar \(U_+\) y \(U_-\) invierte \(H_{\rm lazo}\).
La traza y el conjunto de ganancias extremas son iguales para
\(H_{\rm lazo}\) y \(H_{\rm lazo}^{-1}\). Esos escalares detectan
el efecto del lazo, pero no su orientación. Dos entradas
independientes con sus respuestas vectoriales completas reconstruyen
\(H_{\rm lazo}\); una medida escalar no basta en general.

En el caso conjugado, con ejes de referencia marcados, sean
\[
 G_Q=\frac{E_r(H_{\rm lazo}e_Q)}{E_r(e_Q)},\qquad
 G_P=\frac{E_r(H_{\rm lazo}e_P)}{E_r(e_P)}.
\]
La lectura
\begin{equation}
 \mathcal O=\frac14\log(G_Q/G_P)=4\zeta
 \label{md:comp:lector-orientado}
\end{equation}
recupera la orientación y cambia de signo al invertir el lazo.
Los ejes y la métrica deben transportarse juntos bajo un cambio
de representación. Así se distingue la orientación sin atribuirla
a una traza que la elimina.

\subsection{Información de historia frente a estadística de ocupación}
\label{md:comp:informacion}

Los productos
\[
 U_+U_-U_+^{-1}U_-^{-1}
 \qquad\text{y}\qquad
 U_+U_+^{-1}U_-U_-^{-1}
\]
contienen las mismas cuatro operaciones con las mismas multiplicidades.
El segundo vale \(I\); el primero es \(H_{\rm lazo}\ne I\) para
\(\zeta\ne0\). Por tanto, cualquier estadística que dependa
únicamente de las multiplicidades asigna el mismo resultado a dos
protocolos con acciones diferentes sobre el estado.

La conclusión exacta es que esa estadística no constituye una
descripción suficiente para predecir la respuesta de la composición.
No se refuta Shannon: una distribución sobre secuencias ordenadas
puede representar el orden. Se distingue la información que se ha
conservado en el observable elegido.

\begin{definition}[Información operacional de una historia]
\label{md:comp:def-informacion}
En esta representación, proponemos llamar información operacional de
una historia a la clase de su transporte ordenado distinguible por las
lecturas admitidas. Sea \(\mathcal D\) una colección de detectores
definidos sobre operadores. Puede incluir
\[
 D_{X,L}(H_{\rm lazo})=L(H_{\rm lazo}X),
 \qquad D_{\rm tr}(H_{\rm lazo})=\operatorname{tr}H_{\rm lazo},
\]
con preparación \(X\) y lector lineal \(L\).
Dos historias son equivalentes si
\[
 D(H_{{\rm lazo},1})=D(H_{{\rm lazo},2})
 \qquad\text{para todo }D\in\mathcal D.
\]
\end{definition}

Ampliar la colección puede separar clases antes indistinguibles.
Con respuestas vectoriales completas en una base se reconstruye todo
operador lineal; con sólo traza puede persistir
\(H_{\rm lazo}\sim H_{\rm lazo}^{-1}\).
Así se distingue el detector de una respuesta del estado del detector
espectral del operador.

Esta definición no identifica el transporte con toda la genealogía
TPK: historias distintas pueden producir el mismo operador en una
representación. La información latente en otras hojas, incidencias y
prolongaciones se conserva en sus propios lectores.

La dimensión estadística y la operacional son compatibles.
La tasa de crecimiento combinatorio, o entropía topológica, cuenta
el crecimiento asintótico de configuraciones admisibles; una tasa de
entropía de Shannon requiere además una distribución.
Un mapa de observación determina qué transformaciones pueden
reconstruirse. Ninguna de las pruebas anteriores declara entropía
termodinámica universal nula. El aporte es precisar qué descripción
es suficiente para una tarea de predicción o restitución.

\subsection{Consecuencias y nuevas definiciones propuestas}
\label{md:comp:consecuencias}

Los resultados anteriores reúnen cinco nociones, con sus dominios
y lectores conservados:
\begin{enumerate}
 \item \textbf{Potencial espectral constitutivo.}
 \(\mathscr F\), definido por los momentos \(90/120\), tiene como
 sus dos derivadas el logaritmo de la velocidad normalizada y el
 logaritmo de la impedancia normalizada. Es una definición matemática
 posterior a los canales, sin interpretación energética añadida.

 \item \textbf{Reciprocidad constitutiva diferencial.}
 La igualdad de sensibilidades cruzadas demostrada en la
 sección~\ref{md:comp:reciprocidad} proporciona una restricción
 conjunta verificable, además de los valores.

 \item \textbf{Defecto de cierre relativo.}
 \(\operatorname{tr}H_{\rm lazo}-2\) detecta el lazo no trivial.
 Se acompaña de una lectura orientada, como \(\mathcal O\), cuando
 se requiere distinguir el orden. La traza sola no define toda la memoria.

 \item \textbf{Restitución operacional.}
 Es la recuperación de la acción de la historia desde una colección
 especificada de observables. Se distingue de la restitución de dos
 canales y de la reconstrucción integral del estado.

 \item \textbf{Margen de restitución.}
 La cota \(60a(M_0)\) de la sección~\ref{md:comp:restitucion}
 cuantifica estabilidad absoluta en un dominio, mientras la
 inyectividad sólo afirma unicidad exacta.
\end{enumerate}

Las coordenadas de masa, frecuencia, radios y temperatura del
desarrollo precedente permanecen intactas. Este desarrollo añade
dos niveles: compatibilidad diferencial entre lectores escalares
y memoria de orden en composiciones operatorias.
No se identifica el factor de dilatación de \(H_{\rm lazo}\) con
una masa nueva sin pasar por el lector de masas correspondiente.
